\section{ Nettoyage des textes}

Après avoir réfléchit a commment stocker les romans, il a fallu réfléchir à un moyen de les nettoyer. En effet, les romans contiennent beaucoup de caractères spéciaux qui ne sont pas utiles pour l'analyse, nottament les caractères de mise en forme, pagination, les notes de bas de pages, les chapitres ainsi que prologues (en effet les prologues sont souvent des textes qui ne sont pas écrit par l'auteur du roman et qui ne sont donc pas représentatif de son style d'écriture).

Les romans naturalistes et non naturalistes sont issus de type de données différents. Il y a 3 groupes de format de données : les romans en format \textsl{.epub}, les romans en format \textsl{.pdf} et les romans en format \textsl{.txt}\footnote{ici les tetxes en format \textsl{.txt} sont des en dossier, c'est à dire qu'un fichier \textsl{.txt} correspond a une page du roman, il a donc fallut les regrouper pour obtenir un seul fichier \textsl{.txt} par roman}. 

Pour ce qui est des textes en format \textsl{.epub}, j'ai utilisé le logiciel \textsl{Calibre}\footnote{le lien du lien du locigiciel \href{https://calibre-ebook.com/fr/about}{https://calibre-ebook.com/fr/about}, site consulté le 18 juin 2026} pour les convertir en format \textsl{.txt}. Pour les textes en format \textsl{.pdf}, j'ai utilisé le logiciel \textsl{Adobe Acrobat Pro} pour les convertir en format \textsl{.txt}. Pour les textes en format \textsl{.txt}, j'ai utilisé un script python\footnote{Le script est disponible en annexe} pour regrouper les fichiers en un seul fichier par roman.

La grande majorité de mes romans étaient en format \textsl{.epub} apres leurs conversions en format \textsl{.txt}, j'ai donc utilisé un script python\footnote{Le script est disponible en annexe. Il faut préciser que je n'ai pas utilisé qu'un script mais plusieurs ce sera un script général, il y a eu enormement de modifications interne pour ces derniers notamement pour les chapitres qui sont soit en majuscule, avec chiffre romains ou sans chiffre romains etc. Il a donc fallut faire des ajustemets à chaque fois.} pour nettoyer les textes. Ce script supprime les caractères spéciaux, les notes de bas de pages, les chapitres et les prologues. Il supprime également les lignes vides et les espaces en trop. Le script est basé sur des expressions régulières pour identifier les éléments à supprimer. 

Pour ce qui est du choix des choses a supprimer, j'ai fait le choix de supprimer les chapitres car ils ne sont pas représentatif du style d'écriture de l'auteur. De plus comme nous l'avons vu dans le chapitre 2, j'ai passé un detecteur de motif sur les textes pour identifier les motifs récurrents dans les textes. J'ai donc gardé que le corps de texte pour l'analyse. Ce choix est justifié aussi par le fait d'avoir une noramlisation entres les tetxes\footnote{}.